\section{Practice Questions with Model Answers}
\label{sec:practice}

Worked practice questions spanning all nine question types. Cover the question,
attempt your own answer, then check the model answer.

% =====================================================================
\begin{qbox}[Practice Q1 -- Literacy: confusion matrix (T/F)]
A 4-class action-recognition confusion matrix (rows = true class, columns =
predicted) has rows boxing, handwaving, jogging, running. Row sums are 100.
Diagonals: boxing 98, handwaving 95, jogging 60, running 80. Cell
(boxing,\,handwaving) $=0$; cell (jogging,\,running) $=30$.
Mark each True/False:
(1)~Boxing is never confused with handwaving.
(2)~Jogging is the hardest class to classify.
(3)~Jogging is most often confused with running.
(4)~Running is classified perfectly.

\medskip\noindent\textbf{Model answer.}
(1)~\term{True} -- cell is 0.
(2)~\term{True} -- smallest diagonal (60/100) = lowest recall.
(3)~\term{True} -- its largest off-diagonal cell (30) is the running column.
(4)~\term{False} -- diagonal 80, so 20\% misclassified. Always trace the exact
cell; remember rows are the true classes.
\end{qbox}

% =====================================================================
\begin{qbox}[Practice Q2 -- Algorithms: counterfactual (T/F)]
A loan model rejects an applicant. Mark which statements about \emph{counterfactual
explanations} are correct:
(1)~A good counterfactual changes as few features as possible.
(2)~There is exactly one counterfactual per instance.
(3)~Counterfactuals are model-agnostic and need no access to model internals.
(4)~The counterfactual should be plausible / actionable (e.g.\ not ``reduce age
by 10 years'').

\medskip\noindent\textbf{Model answer.}
Correct: \term{(1), (3), (4)}.
(2) is \term{False} -- many counterfactuals can exist (Rashomon effect); good
methods report close, sparse, plausible ones. (3) holds because counterfactuals
only query input$\to$output behaviour.
\end{qbox}

% =====================================================================
\begin{qbox}[Practice Q3 -- Strategy: 10M traffic-accident records]
A table has 10 million European traffic-accident records: (ID, longitude,
latitude, year). Design an interactive web visualization to explore the
distribution \emph{spatially and temporally}. How would you preprocess?

\medskip\noindent\textbf{Model answer (5-step template).}
\begin{enumerate}[leftmargin=1.4em]
  \item \term{Data type \& task.} Tall (10M rows) spatio-temporal point data;
        task = explore spatial and temporal distribution and their interaction.
  \item \term{Preprocessing / scalability.} The challenge is twofold:
        \emph{perceptual scalability} (10M points overplot any map) and
        \emph{interactive scalability} (raw queries are too slow). Solution:
        \term{spatial binning} (hexbin grid or aggregate to NUTS regions) and
        \term{temporal binning} (by year), with counts pre-aggregated into a
        \term{data cube} over (space-bin $\times$ year) for $O(1)$ slicing.
        Both spatial \emph{and} temporal binning are required.
  \item \term{Visual encoding.} Two views. Spatial: \term{choropleth} (regions)
        or \term{hexbin map}, marks = area cells, channel = colour by accident
        count using a \term{sequential} colour map. Temporal: \term{line chart}
        (or histogram), marks = line/bars, channel = position (x = year,
        y = count).
  \item \term{Interaction.} \term{Linked (coordinated) views}: \term{brush}
        a year range on the temporal chart to update the map, and \term{brush}
        regions on the map to filter the temporal chart (cross-filtering);
        plus zoom/pan and \term{details-on-demand} on hover.
  \item \term{Justify vs.\ alternatives.} Binned map beats raw scatter (no
        overplotting); data cube beats live aggregation (low latency);
        sequential beats rainbow (perceptually ordered). KDE on raw 10M points
        would be too slow / ineffective.
\end{enumerate}
\textit{Deductions to avoid: not naming the scalability challenge; no linked
views; only spatial or only temporal binning; missing marks/channels; rainbow
colour map; no mention of aggregation/data cube.}
\end{qbox}

% =====================================================================
\begin{qbox}[Practice Q4 -- Validation: user study for scatterplot class encodings]
Design a user study to evaluate the effectiveness of visual encodings of
classes/sets in \emph{static} scatterplots.

\medskip\noindent\textbf{Model answer.}
\term{Hypothesis (H1):} colour-coded classes yield faster and more accurate
class-separation judgments than shape-coded classes; \term{H2:} effectiveness
degrades as the number of classes grows.
\term{IV:} class encoding (colour / shape / colour+shape) and \#classes
(3 / 6 / 9).
\term{DV:} task completion time, accuracy/error rate, confidence, preference.
\term{Controlled:} number of points, degree of overlap, point size, display,
task wording, training.
\term{Confounds:} colour-vision deficiency (screen participants), learning/order
effects (counterbalance with a Latin square), fatigue.
\term{Design:} within-subjects, e.g.\ 24 participants; tasks = ``which class is
denser here?'' / ``find the outlier of class~X''; analyze with repeated-measures
ANOVA. Report effect sizes.
\end{qbox}

% =====================================================================
\begin{qbox}[Practice Q5 -- Critique: rainbow choropleth]
A regional-unemployment-rate map uses the rainbow (jet) colour scale. Critique
it and suggest a fix.

\medskip\noindent\textbf{Model answer.}
Problems: (i)~\term{rainbow is not perceptually ordered} -- viewers cannot rank
hues, so high vs.\ low is ambiguous; (ii)~\term{not perceptually uniform} --
equal data steps produce unequal perceived steps (false boundaries at
yellow/cyan); (iii)~\term{not colour-blind safe}; (iv)~mixes hue and lightness,
causing \term{channel interference}; (v)~rate (a magnitude) should use an
ordered channel. Fix: a \term{single-hue (or ColorBrewer) sequential} map
ordered by lightness; consider a \term{diverging} map if comparing against a
national-average reference. Also watch simultaneous contrast between adjacent
regions and the hidden choropleth bias toward large-area regions
(consider rate per capita / a cartogram).
\end{qbox}

% =====================================================================
\begin{qbox}[Practice Q6 -- Alternatives: parallel coordinates vs.\ SPLOM]
The same 8-dimensional dataset is shown as parallel coordinates (PCC) and as a
scatterplot matrix (SPLOM). Compare them.

\medskip\noindent\textbf{Model answer.}
\begin{itemize}[leftmargin=1.4em]
  \item \term{Dimension scalability:} PCC = one axis per dimension (handles many
        dims in one view); SPLOM = $O(d^2)$ panels (64 here), shrinking fast.
  \item \term{Task fit:} SPLOM shows \emph{all pairwise} correlations directly;
        PCC shows \emph{multivariate profiles, clusters, outliers} and
        correlation only between \emph{adjacent} axes.
  \item \term{Clutter:} PCC suffers heavy line overplotting on tall data;
        SPLOM suffers point overplotting per panel.
  \item \term{Ordering:} PCC depends strongly on \term{axis order}
        (non-adjacent relations hidden); SPLOM is order-insensitive.
  \item \term{Interaction:} PCC benefits from brushing + axis reordering; SPLOM
        from linked brushing across panels.
\end{itemize}
Recommendation: SPLOM for pairwise-correlation tasks at moderate $d$; PCC for
profile/cluster/outlier exploration at higher $d$ (with reorderable axes).
\end{qbox}

% =====================================================================
\begin{qbox}[Practice Q7 -- Transfer/sketch: array $\to$ histogram + boxplot]
Given the array $\{4, 5, 5, 6, 6, 6, 7, 8, 20\}$, sketch a histogram and a
boxplot.

\medskip\noindent\textbf{Model answer.}
Sorted $n=9$: median $=6$ (5th value); lower half $\{4,5,5,6\}\Rightarrow$
Q1$=5$; upper half $\{7,8,20\}\Rightarrow$ Q3$=8$; IQR$=3$; upper fence
$=8+1.5\cdot3=12.5\Rightarrow 20$ is an \term{outlier}; whisker max (non-outlier)
$=8$.
\begin{center}
\begin{tikzpicture}[font=\scriptsize, x=0.42cm, y=1cm]
  % ---- boxplot (top) ----
  % axis 0..22 at y=2
  \draw[->,>=Stealth] (-0.5,2) -- (22,2) node[right] {value};
  \foreach \x in {4,5,6,7,8,20}
    \draw (\x,1.92) -- (\x,2.08) node[above=1pt] {\x};
  % whiskers from 4 to 5 and 8 to 8 (max non-outlier = 8)
  \draw[thick] (4,2.6) -- (5,2.6); % lower whisker line
  \draw[thick] (4,2.45) -- (4,2.75); % lower cap
  % box 5..8
  \draw[thick, fill=tuwblue!20] (5,2.3) rectangle (8,2.9);
  % median at 6
  \draw[thick, pitred] (6,2.3) -- (6,2.9);
  % upper whisker (box top = 8 = whisker max, draw small cap)
  \draw[thick] (8,2.45) -- (8,2.75);
  % outlier at 20
  \fill[pitred] (20,2.6) circle (0.16cm);
  \node[pitred, above=2pt] at (20,2.75) {outlier};
  % ---- histogram (bottom), bins width 2 ----
  % bins: [4,6):{4,5,5}=3 ; [6,8):{6,6,6,7}=4 ; [8,10):{8}=1 ; [20,22):{20}=1
  \draw[->,>=Stealth] (-0.5,0) -- (22,0) node[right] {value};
  \draw[->,>=Stealth] (-0.5,0) -- (-0.5,1.3) node[above] {count};
  \fill[keygreen!50] (4,0) rectangle (6,0.9);
  \fill[keygreen!50] (6,0) rectangle (8,1.2);
  \fill[keygreen!50] (8,0) rectangle (10,0.3);
  \fill[keygreen!50] (20,0) rectangle (22,0.3);
  \foreach \x in {4,6,8,10,20,22}
    \draw (\x,-0.06) -- (\x,0.06) node[below=2pt] {\x};
\end{tikzpicture}
\end{center}
\sketchnote{Boxplot: box from 5 to 8, median line at 6, lower whisker to 4,
upper whisker to 8, a separate dot at 20 marked as outlier. Histogram (bins of
width 2): tall bar around 5--6, smaller bars at 4 and 7--8, an isolated low bar
near 20 -- right-skewed.}
\end{qbox}

% =====================================================================
\begin{qbox}[Practice Q8 -- Mapping: scatterplot point $\to$ PCC polyline]
A scatterplot of (weight, mpg) highlights one outlier car: high weight, low mpg.
Mark the same car in a parallel-coordinates plot whose axes are, left to right,
mpg, weight, horsepower.

\medskip\noindent\textbf{Model answer.}
Recover its values: low on mpg, high on weight (horsepower unknown from the
scatter, but heavy cars are typically high). The car's \term{polyline} is the
one that is \term{low on the mpg axis and high on the weight axis} (and, for a
heavy car, high on horsepower). Trace the single line crossing those positions
and highlight it.
\begin{center}
\begin{tikzpicture}[font=\scriptsize]
  % three vertical axes at x=0,2.5,5 ; y from 0 (low) to 4 (high)
  \foreach \x/\lab in {0/mpg, 2.5/weight, 5/horsepower}{
    \draw[thick] (\x,0) -- (\x,4);
    \node[below] at (\x,0) {\lab};
    \node[above=2pt] at (\x,4) {high};
    \node[below=10pt] at (\x,0) {low};
  }
  % faded background lines
  \draw[gray!55] (0,3.0) -- (2.5,1.4) -- (5,2.2);
  \draw[gray!55] (0,2.5) -- (2.5,2.0) -- (5,1.6);
  \draw[gray!55] (0,3.4) -- (2.5,1.0) -- (5,1.2);
  \draw[gray!55] (0,2.0) -- (2.5,2.6) -- (5,2.8);
  % highlighted outlier: low mpg, high weight, high horsepower
  \draw[ultra thick, pitred] (0,0.5) -- (2.5,3.7) -- (5,3.5);
  \fill[pitred] (0,0.5) circle (1.6pt) (2.5,3.7) circle (1.6pt) (5,3.5) circle (1.6pt);
\end{tikzpicture}
\end{center}
\sketchnote{Three vertical axes; one bold polyline starting low on mpg, rising
to the top of the weight axis, staying high on horsepower; other lines faded.}
\end{qbox}

% =====================================================================
\begin{qbox}[Practice Q9 -- MCQ: t-SNE properties]
Which statements about t-SNE are correct?
(a)~It is a linear projection.
(b)~Perplexity balances attention to local vs.\ global structure.
(c)~Distances between well-separated clusters are quantitatively meaningful.
(d)~Different random seeds can give different layouts.

\medskip\noindent\textbf{Model answer.} Correct: \term{(b), (d)}.
(a) is false (nonlinear); (c) is false (between-cluster distances and cluster
sizes are \emph{not} reliable). t-SNE is stochastic, so set a seed for
reproducibility.
\end{qbox}

% =====================================================================
\begin{qbox}[Practice Q10 -- MCQ: PCA vs.\ UMAP]
Pick the correct pairings.
(a)~PCA axes are interpretable directions of maximum variance.
(b)~UMAP is deterministic.
(c)~UMAP's \texttt{n\_neighbors} trades local detail for global structure.
(d)~PCA can capture nonlinear manifolds.

\medskip\noindent\textbf{Model answer.} Correct: \term{(a), (c)}.
(b) false -- UMAP is stochastic. (d) false -- PCA is linear; use t-SNE/UMAP/MDS
or kernel PCA for nonlinear structure.
\end{qbox}

% =====================================================================
\begin{qbox}[Practice Q11 -- MCQ: force-directed layout]
Which hold for force-directed graph layout?
(a)~It is deterministic regardless of initialization.
(b)~It models edges as springs and nodes as mutually repelling charges.
(c)~Naive computation is $O(n^2)$ per iteration; Barnes--Hut approximates it.
(d)~It guarantees the global minimum-energy layout.

\medskip\noindent\textbf{Model answer.} Correct: \term{(b), (c)}.
(a) false -- the result depends on the random initial placement. (d) false --
it can settle in a local minimum.
\end{qbox}

% =====================================================================
\begin{qbox}[Practice Q12 -- List: overplotting fixes with no preprocessing]
List simple solutions to overplotting that require \emph{no} data preprocessing.

\medskip\noindent\textbf{Model answer.}
\term{Transparency (alpha)}, \term{smaller marks}, \term{open glyphs},
\term{jitter}, \term{on-the-fly subsampling}. (Hexbin, KDE, aggregation, and
DR are effective but \emph{do} require preprocessing, so they do not count.)
\end{qbox}

% =====================================================================
\begin{qbox}[Practice Q13 -- Algorithms: sketch 1D PCA]
Sketch the 1D PCA projection of an elongated 2D cloud tilted at about
$30^{\circ}$.

\medskip\noindent\textbf{Model answer.}
\begin{center}
\begin{tikzpicture}[font=\scriptsize, rotate=30]
  % elongated cloud tilted 30 deg (whole picture rotated)
  \draw[tuwblue, dashed] (0,0) ellipse (2.6cm and 0.8cm);
  % PC1 arrow along long axis (x), PC2 short axis (y)
  \draw[->,>=Stealth, ultra thick, keygreen] (-2.9,0) -- (2.9,0) node[right] {PC1};
  \draw[->,>=Stealth, thick, gray] (0,-1.1) -- (0,1.1) node[above] {PC2 (discarded)};
  % data points and their perpendicular projections (feet) onto PC1
  \foreach \px/\py in {-2.0/0.4, -1.2/-0.3, -0.4/0.5, 0.3/-0.2, 1.0/0.4, 1.8/-0.35, -1.6/-0.5, 0.7/0.5}{
    \fill[pitred] (\px,\py) circle (1.4pt);
    \draw[gray!60, thin] (\px,\py) -- (\px,0);
    \fill[keygreen] (\px,0) circle (1.0pt);
  }
\end{tikzpicture}
\end{center}
\sketchnote{Draw the tilted elliptical cloud. Draw PC1 as an arrow along its
long axis (max variance). Project each point perpendicularly onto PC1; the feet
of those perpendiculars are the 1D coordinates. PC2 (orthogonal, short axis) is
discarded.}
PCA is linear and deterministic; PC1 captures the most variance.
\end{qbox}

% =====================================================================
\begin{qbox}[Practice Q14 -- List: model-agnostic interpretability methods]
Name model-agnostic methods to interpret an ML model.

\medskip\noindent\textbf{Model answer.}
\term{PDP} (average marginal effect), \term{ICE} (per-instance curves),
\term{permutation feature importance}, \term{LIME} (local surrogate),
\term{SHAP} (Shapley attribution), \term{counterfactual explanations},
\term{ALE plots}, and global \term{surrogate models}. All treat the model as a
black box, querying only inputs$\to$outputs.
\end{qbox}
